\documentclass{amsart}
\usepackage{fullpage,graphicx,subfigure,mathpazo,color}
\usepackage{amsmath,amscd,tikz,mathrsfs}
\usepackage[normalem]{ulem}
\usepackage{amsmath}
\usepackage{epstopdf}
\usepackage{tikz}
\usepackage{setspace}%ʹÓÃŒäŸàºê°ü
\usepackage{hyperref}
%\usepackage[backref]{hyperref}
%\usepackage[ctex]{ulem}
\newcommand{\ii}{\mathrm{i}}
\newcommand{\ee}{\mathrm{e}}
\newcommand{\dd}{\mathrm{d}}
\newcommand{\T}{\mathrm{T}}
\newcommand{\lm}[1]{{\color{red} #1}}
\newcommand{\xe}[1]{{\color{blue} #1}}

\newcommand{\cn}{\mathrm{cn}}
\newcommand{\sn}{\mathrm{sn}}
\newcommand{\dn}{\mathrm{dn}}


\newtheorem{rhp}{Riemann-Hilbert Problem}
\newtheorem{thm}{Theorem}
\newtheorem{lem}{Lemma}
\newtheorem{prop}{Proposition}
\newtheorem{corollary}{Corollary}
\newtheorem{definition}{Definition}
\newtheorem{conjecture}{Conjecture}
\newtheorem{rem}{Remark}
\newtheorem{example}{Example}

\usepackage{graphicx}      % insert graphic
%\usepackage{subcaption}
\usepackage{titlesec}

\newcounter{week}
\setcounter{week}{0}

% Week macro: creates a section for the week, adds to TOC, and resets local numbering if desired
\newcommand{\Week}[2]{%
  \refstepcounter{week}%
  \subsection*{Week \theweek: #1 \hfill {\small #2}}%
  \addcontentsline{toc}{section}{Week \theweek: #1 --- #2}%
  % \setcounter{theorem}{0} % optional: reset theorem numbering each week
}


% Smaller helper macros commonly useful in notes
\newcommand{\obj}[1]{\paragraph{Objectives:} #1}
\newcommand{\reading}[1]{\paragraph{Reading:} #1}
\newcommand{\summary}{\paragraph{Lecture summary:}}
\newcommand{\defs}{\paragraph{Definitions:}}
\newcommand{\thms}{\paragraph{Theorems \& Propositions:}}
\newcommand{\exs}{\paragraph{Examples:}}
\newcommand{\exer}{\paragraph{Exercises:}}
\newcommand{\notes}{\paragraph{Notes and remarks:}}
\newcommand{\refs}{\paragraph{References:}}



\titleformat{\section}{\centering\LARGE\bfseries}{\thesection}{1em}{}
\titleformat{\subsection}{\Large\bfseries}{\thesubsection}{1em}{}
\begin{document}


\title{Lecture Notes on Integrable Systems}

% \author{Liming Ling}
% \address{School of Mathematics, South China University of Technology, Guangzhou, China 510641}
% \email{linglm@scut.edu.cn}
\author{Zixuan Deng}
\address{School of Mathematics, South China University of Technology, Guangzhou}
\email{2627441954@qq.com}
\date{\today}

\maketitle


% ---------- Table of contents for a notes file ----------
% \begin{center}
%   \textbf{\large Course Notes (weekly)}\\[4pt]
%   \textit{--- Table of Contents (auto-generated) ---}
% \end{center}
% \smallskip
% \tableofcontents
% \bigskip
% \hrule
% \bigskip


\Week{Algebro-Geometric Solutions for the mKdV Equation (1)}{2025-09-16}

\obj{Derive the mKdV equation from the associated spectral problem and construct the AKNS hierarchy by explicit recursion. } 

% \summary{Starting from the Lax pair, we expand the formal solution, obtain recursive relations for the coefficients, and use them to derive the evolution equations. This yields the mKdV equation as a member of the AKNS hierarchy, with all steps carried out explicitly.} 


The modified Korteweg--de Vries (mKdV) equation is given by
\begin{equation}\label{eq:mkdv}
u_t + 6u^2 u_x + u_{xxx} = 0,
\qquad u(x,t) \in \mathbb{R}, \quad (x,t) \in \mathbb{R}^2.
\end{equation}

A standard method for establishing integrability is via a Lax pair. We introduce a matrix function $\Phi(x,t)$ satisfying
\begin{equation*}
\Phi_x = U \Phi, 
\qquad 
\Phi_t = V \Phi,
\end{equation*}
where $U$ and $V$ are $2 \times 2$ matrices depending on $u$ and its derivatives. The compatibility condition
\begin{equation*}
\Phi_{xt} = \Phi_{tx}
\quad \Longleftrightarrow \quad
U_t - V_x + [U,V] = 0,
\end{equation*}
with $[U,V] = UV - VU$, is the zero-curvature equation.

Following the AKNS formalism, the spatial part of the Lax pair is taken as
\begin{equation}\label{eq:AKNS-x}
\Phi_x = U \Phi,
\qquad
U = -\mathrm{i} \lambda \sigma_3 + Q,
\end{equation}
where
\begin{equation}\label{eq:Q-sigma3}
Q = \begin{pmatrix} 0 & u \\ -u & 0 \end{pmatrix},
\qquad
\sigma_3 = \operatorname{diag}(1,-1).
\end{equation}

We introduce the transformation
\begin{equation}\label{eq:Phi-gauge}
\Phi(x;\lambda) = m(x;\lambda) e^{-\mathrm{i}\lambda x \sigma_3}, 
\qquad 
m(x;\lambda) \to \mathbb{I} \quad \text{as } \lambda \to \infty.
\end{equation}
Substituting into the Lax equation gives
\begin{equation}\label{eq:lax}
\Phi_x \Phi^{-1} 
= m_x m^{-1} - \mathrm{i}\lambda m\sigma_3 m^{-1}
= -\mathrm{i}\lambda \sigma_3 + Q.
\end{equation}
Define
\begin{equation}\label{eq:Psi-def}
\Psi = m \sigma_3 m^{-1}.
\end{equation}
Differentiating $\Psi$ and using \eqref{eq:lax}, we compute
\begin{align*}
\Psi_x 
&= m_x \sigma_3 m^{-1} - m \sigma_3 m^{-1} m_x m^{-1} \\
&= [m_x m^{-1}, m \sigma_3 m^{-1}] \\
&= \big[m_x m^{-1} - \mathrm{i}\lambda m\sigma_3 m^{-1}, m \sigma_3 m^{-1}\big] \\
&= [U,\Psi].
\end{align*}
Thus $\Psi$ satisfies
\begin{equation}\label{eq:zero-curvature}
\Psi_x = [U,\Psi],
\end{equation}
which is the stationary zero-curvature condition.

The matrix $m(x;\lambda)$ admits the asymptotic expansion
\begin{equation}\label{eq:m-expansion}
m(x;\lambda) = \mathbb{I} + \lambda^{-1} m_1 + \lambda^{-2} m_2 + \cdots.
\end{equation}
From \eqref{eq:Psi-def} and \eqref{eq:m-expansion} we obtain the formal expansion
\begin{equation}\label{eq:Psi-expansion}
\Psi = \Psi_0 + \Psi_1 \lambda^{-1} + \Psi_2 \lambda^{-2} + \cdots,
\qquad 
\Psi_0 = \sigma_3.
\end{equation}
Substituting this into \eqref{eq:zero-curvature} yields
\begin{align*}
\sum_{j=1}^{\infty} \Psi_{j,x} \lambda^{-j} 
&= - \mathrm{i}\lambda \Big[ \sigma_3, \sum_{j=1}^{\infty} \Psi_j \lambda^{-j} \Big]
+ \Big[ Q, \sum_{j=0}^\infty \Psi_j \lambda^{-j} \Big] \\
&= - \mathrm{i}\sum_{j=0}^{\infty} \Big[ \sigma_3,  \Psi_{j+1}  \Big] \lambda^{-j}
+ \sum_{j=0}^\infty \Big[ Q,  \Psi_j \Big] \lambda^{-j}.
\end{align*}

At order $O(1)$ we find
\begin{equation*}
- \mathrm{i} [\sigma_3, \Psi_1] + [Q, \sigma_3] = 0 
\quad \Longleftrightarrow \quad [\mathrm{i}\Psi_1 + Q, \sigma_3] = 0,
\end{equation*}
which implies
\begin{equation*}
\Psi_1^{\mathrm{off}} = \mathrm{i} Q.
\end{equation*}

For general $j \geq 1$, the coefficient of $\lambda^{-j}$ gives
\begin{equation}\label{eq:Psi-j-recursion}
\Psi_{j,x} 
= - \mathrm{i}[\sigma_3,\Psi_{j+1}] + [Q,\Psi_j]
= -2\mathrm{i}\sigma_3 \Psi_{j+1}^{\mathrm{off}} + [Q,\Psi_j].
\end{equation}
Hence
\begin{equation}\label{eq:Psi-j-off-recursion}
\Psi_{j,x}^{\mathrm{off}} 
= -2\mathrm{i} \sigma_3 \Psi_{j+1}^{\mathrm{off}} + [Q,\Psi_j^{\mathrm{diag}}]
\quad \Longleftrightarrow \quad
\Psi_{j+1}^{\mathrm{off}} 
= \tfrac{\mathrm{i}}{2}\sigma_3\big(\Psi_{j,x}^{\mathrm{off}} - [Q,\Psi_j^{\mathrm{diag}}]\big),
\end{equation}
and
\begin{equation}\label{eq:Psi-j-diag-recursion}
\Psi_{j,x}^{\mathrm{diag}} = [Q,\Psi_j^{\mathrm{off}}]
\quad \Longleftrightarrow \quad
\Psi_j^{\mathrm{diag}} = \partial_x^{-1}\big([Q,\Psi_j^{\mathrm{off}}]\big).
\end{equation}

Thus, the recursion relation \eqref{eq:Psi-j-off-recursion} provides a direct way to compute the
off-diagonal coefficients algebraically in terms of lower-order terms, while the diagonal recursion
\eqref{eq:Psi-j-diag-recursion} requires performing an integration. In practice,
this makes the off-diagonal part more convenient to obtain explicitly, 
whereas the diagonal part is generally harder to compute in closed form.

By \eqref{eq:Psi-def}, the additional constraint $\Psi^2 = \mathbb{I}$ provides algebraic relations among the coefficients. Substituting the expansion \eqref{eq:Psi-expansion} yields
\begin{equation*}
\mathbb{I} = \left( \sum_{j=0}^\infty \Psi_j \lambda^{-j} \right)^2
= \sum_{j=0}^\infty \left( \sum_{k=0}^j \Psi_k \Psi_{j-k} \right) \lambda^{-j},
\end{equation*}
so that
\begin{equation}\label{eq:Psi-sq-coeff}
\sum_{k=0}^j \Psi_k \Psi_{j-k} = 0, 
\qquad j \geq 1.
\end{equation}
Separating diagonal and off-diagonal parts, for $j \geq 1$ we obtain
\begin{equation}\label{eq:Psi-sq-diag}
    \sum_{k=0}^j \big( \Psi_k^{\mathrm{diag}} \Psi_{j-k}^{\mathrm{diag}} + \Psi_k^{\mathrm{off}} \Psi_{j-k}^{\mathrm{off}} \big) = 0,
\end{equation}
\begin{equation}\label{eq:Psi-sq-off-diag}
\sum_{k=0}^j \big( \Psi_k^{\mathrm{off}} \Psi_{j-k}^{\mathrm{diag}} + \Psi_k^{\mathrm{diag}} \Psi_{j-k}^{\mathrm{off}} \big) = 0.
\end{equation}

Since $\Psi_0 = \sigma_3$, the off-diagonal constraint \eqref{eq:Psi-sq-off-diag} for $j \geq 2$ takes the form
\begin{equation*}
\sigma_3 \Psi_j^{\mathrm{off}} + \Psi_j^{\mathrm{off}} \sigma_3 
+ \sum_{k=1}^{j-1} \left( \Psi_k^{\mathrm{off}} \Psi_{j-k}^{\mathrm{diag}} + \Psi_k^{\mathrm{diag}} \Psi_{j-k}^{\mathrm{off}} \right) = 0.
\end{equation*}
Noting that $\sigma_3 \Psi_j^{\mathrm{off}} + \Psi_j^{\mathrm{off}} \sigma_3 = 0$, this simplifies to
\begin{equation}\label{eq:Psi-sq-off-diag-reduced}
\sum_{k=1}^{j-1} \left( \Psi_k^{\mathrm{off}} \Psi_{j-k}^{\mathrm{diag}} + \Psi_k^{\mathrm{diag}} \Psi_{j-k}^{\mathrm{off}} \right) = 0.
\end{equation}

It should be noted that from the off-diagonal constraint \eqref{eq:Psi-sq-off-diag} one can only deduce 
the reduced form \eqref{eq:Psi-sq-off-diag-reduced}, which merely expresses a compatibility condition 
between lower-order coefficients. In particular, this relation does not provide a practical means 
to compute the diagonal parts explicitly.

For the diagonal part \eqref{eq:Psi-sq-diag} we have
\begin{equation*}
\sigma_3 \Psi_j^{\mathrm{diag}} + \Psi_j^{\mathrm{diag}} \sigma_3 
+ \sum_{k=1}^{j-1} \big( \Psi_k^{\mathrm{diag}} \Psi_{j-k}^{\mathrm{diag}} + \Psi_k^{\mathrm{off}} \Psi_{j-k}^{\mathrm{off}} \big) = 0,
\qquad j \geq 2.
\end{equation*}
Since $\sigma_3 \Psi_j^{\mathrm{diag}} = \Psi_j^{\mathrm{diag}} \sigma_3$, this simplifies to
\begin{equation}\label{eq:Psi-sq-diag-reduced}
\Psi_j^{\mathrm{diag}} 
= -\tfrac{1}{2} \sigma_3 \sum_{k=1}^{j-1} 
\big( \Psi_k^{\mathrm{diag}} \Psi_{j-k}^{\mathrm{diag}} + \Psi_k^{\mathrm{off}} \Psi_{j-k}^{\mathrm{off}} \big),
\qquad j \geq 2.
\end{equation}
We will use this reduced relation \eqref{eq:Psi-sq-diag-reduced} to compute the diagonal part $\Psi_j^{\mathrm{diag}}$ recursively at each order. This provides an explicit algebraic formula for the diagonal coefficients in terms of lower-order diagonal and off-diagonal terms.

For $j=1$, the condition \eqref{eq:Psi-sq-coeff} reads
\begin{equation*}
\sigma_3 \Psi_1 + \Psi_1 \sigma_3 = 2\sigma_3 \Psi_1^{\mathrm{diag}} = 0,
\end{equation*}
so that $\Psi_1 = \mathrm{i}Q$.

Using \eqref{eq:Psi-j-off-recursion} and \eqref{eq:Psi-sq-diag-reduced} recursively, we find
\begin{equation*}
\Psi_2^{\mathrm{off}} = -\tfrac{1}{2} \sigma_3 Q_x,
\qquad
\Psi_2^{\mathrm{diag}} = \tfrac{1}{2} \sigma_3 Q^2,
\end{equation*}
and hence
\begin{equation*}
\Psi_2 = \tfrac{1}{2} \sigma_3 (Q^2 - Q_x).
\end{equation*}

For the next coefficient, we compute
\begin{align*}
\Psi_3^{\mathrm{off}} 
&= \tfrac{\mathrm{i}}{2} \sigma_3 \big( \Psi_{2,x}^{\mathrm{off}} - [Q,\Psi_2^{\mathrm{diag}}] \big) \\
&= \tfrac{\mathrm{i}}{2} \sigma_3 \left( -\tfrac{1}{2} \sigma_3 Q_{xx} - [Q, \tfrac{1}{2} \sigma_3 Q^2] \right) \\
&= \tfrac{\mathrm{i}}{2} \sigma_3 \left( -\tfrac{1}{2} \sigma_3 Q_{xx} - \tfrac{1}{2}(Q\sigma_3 Q^2 - \sigma_3 Q^2 Q) \right) \\
&= -\tfrac{\mathrm{i}}{4} Q_{xx} + \tfrac{\mathrm{i}}{2} Q^3.
\end{align*}
For the diagonal part,
\begin{align*}
\Psi_3^{\mathrm{diag}} 
&= -\tfrac{1}{2} \sigma_3 \Big( \Psi_1^{\mathrm{off}}\Psi_2^{\mathrm{off}} + \Psi_2^{\mathrm{off}}\Psi_1^{\mathrm{off}} 
+ \Psi_1^{\mathrm{diag}}\Psi_2^{\mathrm{diag}} + \Psi_2^{\mathrm{diag}}\Psi_1^{\mathrm{diag}} \Big) \\
&= -\tfrac{1}{2} \sigma_3 \Big( (\mathrm{i}Q)(-\tfrac{1}{2}\sigma_3 Q_x) + (-\tfrac{1}{2}\sigma_3 Q_x)(\mathrm{i}Q) \Big) \\
&= -\tfrac{1}{2} \sigma_3 \left( -\tfrac{\mathrm{i}}{2}(Q\sigma_3 Q_x + \sigma_3 Q_x Q) \right) \\
&= \tfrac{\mathrm{i}}{4}\big(Q_x Q - Q Q_x\big) = 0.
\end{align*}

Proceeding to $\Psi_4$, the off-diagonal recursion \eqref{eq:Psi-j-off-recursion} with $\Psi_3^{\mathrm{diag}}=0$ gives
\begin{align*}
\Psi_4^{\mathrm{off}}
&= \tfrac{\mathrm{i}}{2}\sigma_3\big(\Psi_{3,x}^{\mathrm{off}} - [Q,\Psi_3^{\mathrm{diag}}]\big)
= \tfrac{\mathrm{i}}{2}\sigma_3 \Psi_{3,x}^{\mathrm{off}} \\
&= \tfrac{\mathrm{i}}{2}\sigma_3 \partial_x\!\big(-\tfrac{\mathrm{i}}{4} Q_{xx} + \tfrac{\mathrm{i}}{2} Q^3\big) \\
&= \tfrac{\mathrm{i}}{2}\sigma_3 \big(-\tfrac{\mathrm{i}}{4} Q_{xxx} + \tfrac{\mathrm{i}}{2} (Q^3)_x\big) \\
&= \tfrac{1}{8}\sigma_3 Q_{xxx} - \tfrac{1}{4}\sigma_3 (Q^3)_x.
\end{align*}
Since $Q_x Q = Q Q_x$ for the present $Q$ given in \eqref{eq:Q-sigma3}, we have
\[
(Q^3)_x = Q_x Q^2 + Q Q_x Q + Q^2 Q_x = 3 Q_x Q^2,
\]
and therefore
\begin{equation*}
\Psi_4^{\mathrm{off}} = \tfrac{1}{8}\sigma_3 Q_{xxx} - \tfrac{3}{4}\sigma_3 Q_x Q^2.
\end{equation*}
For the diagonal part, only the nonvanishing combinations contribute; terms involving $\Psi_1^{\mathrm{diag}}$ or $\Psi_3^{\mathrm{diag}}$ are zero and omitted. Hence
\begin{align*}
\Psi_4^{\mathrm{diag}}
&= -\tfrac{1}{2}\sigma_3\Big(\Psi_1^{\mathrm{off}}\Psi_3^{\mathrm{off}} + \Psi_3^{\mathrm{off}}\Psi_1^{\mathrm{off}}
+ \Psi_2^{\mathrm{off}}\Psi_2^{\mathrm{off}} + \Psi_2^{\mathrm{diag}}\Psi_2^{\mathrm{diag}}\Big) \\
&= -\tfrac{1}{2}\sigma_3\Big((\mathrm{i}Q)\big(-\tfrac{\mathrm{i}}{4}Q_{xx}+\tfrac{\mathrm{i}}{2}Q^3\big)
+ \big(-\tfrac{\mathrm{i}}{4}Q_{xx}+\tfrac{\mathrm{i}}{2}Q^3\big)(\mathrm{i}Q) \\
&\qquad\qquad\quad + \big(-\tfrac{1}{2}\sigma_3 Q_x\big)^2 + \big(\tfrac{1}{2}\sigma_3 Q^2\big)^2\Big) \\
&= -\tfrac{1}{8}\sigma_3\big(Q Q_{xx} + Q_{xx} Q - 3Q^4 - Q_x^2\big).
\end{align*}

Since $(-\mathrm{i}\lambda \Psi)_{+} = -\mathrm{i}\lambda \sigma_3 + Q$, we obtain
\begin{equation*}
\Phi_x \Phi^{-1} = U = (-\mathrm{i}\lambda \Psi)_{+}.
\end{equation*}
Here the subscript $(\cdot)_{+}$ denotes the projection onto the nonnegative powers of $\lambda$ in the Laurent expansion.  
This motivates the definition of the AKNS hierarchy by prescribing the time evolution
\begin{equation}\label{eq:AKNS-hierarchy}
\Phi_{t_n} = \big(-\mathrm{i}\lambda^n \Psi\big)_{+}\Phi \equiv V_n \Phi.
\end{equation}

The first few members of the hierarchy are given by
\begin{align*}
n=2 &\quad \Rightarrow \quad \text{nonlinear Schrödinger (NLS) flow}, \\
n=3 &\quad \Rightarrow \quad \text{complex mKdV flow}, \quad \Rightarrow \quad \text{mKdV flow}.
\end{align*}
In terms of the expansion coefficients $\Psi_j$, the operators $V_n$ take the form
\begin{align*}
V_2 &= -\mathrm{i}\big(\lambda^2 \Psi_0 + \lambda \Psi_1 + \Psi_2\big), \\
V_3 &= -\mathrm{i}\big(\lambda^3 \Psi_0 + \lambda^2 \Psi_1 + \lambda \Psi_2 + \Psi_3\big).
\end{align*}

The zero-curvature condition
\begin{equation*}
U_t - V_x + [U,V] = 0,
\end{equation*}
with $V_{\mathrm{mKdV}} = 4 V_3$, yields at order $O(1)$
\begin{equation*}
Q_t + 4\mathrm{i} \Psi_{3,x} + [Q,-4\mathrm{i}\Psi_3] = 0.
\end{equation*}
Taking the off-diagonal part gives
\begin{equation*}
Q_t + 4\mathrm{i}\Psi_{3,x}^{\mathrm{off}} = 0,
\end{equation*}
that is,
\begin{align}\label{eq:matrix-mkdv}    
Q_t + 4\mathrm{i}\Big(-\tfrac{\mathrm{i}}{4}Q_{xx}+\tfrac{\mathrm{i}}{2}Q^3\Big)_x
&= Q_t + Q_{xxx} - 2(Q^3)_x = 0.
\end{align}
Since
\begin{equation*}
Q = \begin{pmatrix} 0 & u \\ -u & 0 \end{pmatrix}, 
\qquad Q^2 = -u^2 \mathbb{I},
\end{equation*}
this reduces to the scalar equation
\begin{equation*}
u_t + u_{xxx} + 2(u^3)_x = 0,
\end{equation*}
which is precisely the mKdV equation \eqref{eq:mkdv}.

Using the definitions \eqref{eq:Phi-gauge} and  \eqref{eq:Psi-def} we have
\begin{equation}\label{eq:Psi-def-new}
\Psi = \Phi \sigma_3 \Phi^{-1}.
\end{equation}
Differentiating $\Psi$ with respect to $t_n$ and using \eqref{eq:AKNS-hierarchy} gives
\begin{align*}
\Psi_{t_n}
&= \Phi_{t_n}\sigma_3 \Phi^{-1} - \Phi\sigma_3 \Phi^{-1} \Phi_{t_n} \Phi^{-1} \\
&= [\Phi_{t_n} \Phi^{-1},\, \Psi] = [V_n,\Psi].
\end{align*}
Thus
\begin{equation}\label{eq:Psi-t-comm}
\Psi_{t_n} = [V_n,\Psi].
\end{equation}

Next we simplify the zero-curvature condition for the pair $(U,V_n)$.
The zero-curvature equation is
\[
U_{t_n} - (V_n)_x + [U,V_n]=0,
\]
and since $U=-\mathrm{i}\lambda\sigma_3+Q$ we have $U_{t_n}=Q_{t_n}$.  
With the normalization
\[
V_n=\alpha_n\big(-\mathrm{i}\lambda^n\Psi\big)_{+}
\]
we obtain
\begin{equation*}
Q_{t_n}=(V_n)_x - [U,V_n]
= \alpha_n\big(-\mathrm{i}\lambda^n\Psi_x\big)_{+} - [U,\alpha_n\big(-\mathrm{i}\lambda^n\Psi\big)_{+}].
\end{equation*}
Pulling out constants and using $\Psi_x=[U,\Psi]$ gives
\begin{equation}\label{eq:interm-refined}
Q_{t_n} = -\mathrm{i}\alpha_n\Big(\big(\lambda^n[U,\Psi]\big)_{+} - [U,\big(\lambda^n\Psi\big)_{+}]\Big).
\end{equation}

Write $U=-\mathrm{i}\lambda\sigma_3+Q$ and split the first term in \eqref{eq:interm-refined}:
\[
\lambda^n[U,\Psi] = -\mathrm{i}\lambda^{n+1}[\sigma_3,\Psi] + \lambda^n[Q,\Psi].
\]
Because $Q$ does not depend on $\lambda$ we have
\[
\big(\lambda^n[Q,\Psi]\big)_{+} = [Q,\big(\lambda^n\Psi\big)_{+}],
\]
so the $\!Q$--commutator contributions cancel in \eqref{eq:interm-refined}. The remaining terms therefore come only from the $\sigma_3$--commutator:
\begin{align*}
Q_{t_n}
&= -\mathrm{i}\alpha_n\Big(-\mathrm{i}\big(\lambda^{n+1}[\sigma_3,\Psi]\big)_{+} + \mathrm{i}\,\lambda[\sigma_3,\big(\lambda^n\Psi\big)_{+}]\Big)\\
&= \alpha_n\Big(\lambda[\sigma_3,\big(\lambda^n\Psi\big)_{+}] - \big(\lambda^{n+1}[\sigma_3,\Psi]\big)_{+}\Big).
\end{align*}

To simplify the difference of finite sums, use the expansions
\[
\Psi=\sum_{j\ge0}\Psi_j\lambda^{-j},\qquad
(\lambda^n\Psi)_{+}=\sum_{j=0}^n\lambda^{\,n-j}\Psi_j.
\]
A direct comparison of coefficients shows
\begin{align*}
\lambda[\sigma_3,(\lambda^n\Psi)_{+}] - \big(\lambda^{n+1}[\sigma_3,\Psi]\big)_{+}
&= \sum_{j=0}^n \lambda^{\,n+1-j}[\sigma_3,\Psi_j] - \sum_{j=0}^{n+1} \lambda^{\,n+1-j}[\sigma_3,\Psi_j] \\
&= -\lambda^{\,n+1-(n+1)}[\sigma_3,\Psi_{n+1}] \\
&= -[\sigma_3,\Psi_{n+1}].
\end{align*}
Hence the evolution of $Q$ is
\begin{equation}\label{eq:Q-t-final-refined}
Q_{t_n} = -\alpha_n\,[\sigma_3,\Psi_{n+1}].
\end{equation}
Since diagonal parts commute with $\sigma_3$ and the commutator doubles the off-diagonal part, we may equivalently write
\begin{equation}\label{eq:Q-t-off}
Q_{t_n} = -2\alpha_n\,\sigma_3\,\Psi_{n+1}^{\mathrm{off}}.
\end{equation}





% \bibliographystyle{siam}
% \bibliography{CFL-pattern}





\end{document}

